\section{The Old Preacher and \work{Against Celsus}}

Persecution shadowed the Caesarean years as it had the Alexandrian. Under
Maximin Thrax (235--238), who ordered ``that only the rulers of the churches
should be put to death, as responsible for the Gospel teaching,'' Origen's
patron Ambrose and the Caesarean presbyter Protoctetus were seized, and
``thereupon Origen composed his work \work{On Martyrdom},'' the
\work{Exhortation to Martyrdom}, dedicating it to them (\work{Church
History} 6.28). The son of Leonides, now past fifty, wrote for his friends
the treatise he had once, at sixteen, compressed into a single line to his
father.

Age changed his practice in one respect Eusebius thought worth recording:
``being, as they say, over sixty years old, and having gained great facility
by his long practice,'' Origen finally ``permitted his public discourses to
be taken down by stenographers, a thing which he had never before allowed''
(\work{Church History} 6.36.1). Most of the surviving homilies---the largest
body of Christian preaching before the fourth century---descend from this
late permission. The preacher they show is far from the speculative
metaphysician of controversy: a pastor working book by book through
Scripture, impatient of long services and inattentive congregations,
relentlessly moral and Christ-centered.

To the same last decade belongs the greatest ancient defense of
Christianity. ``He also at this time composed a work of eight books in
answer to that entitled \work{True Discourse}, which had been written
against us by Celsus the Epicurean'' (\work{Church History} 6.36.2)---in
fact a Platonist critic, writing about seventy years earlier, whose book
survives only in the massive extracts Origen quotes in order to refute.
Ambrose had sent the book and asked for the reply. Origen's preface is a
reluctant masterpiece: Jesus, ``when false witness was borne against Him,
remained silent,'' and even now ``continues silent before these things, and
makes no audible answer, but places His defense in the lives of His genuine
disciples, which are a pre-eminent testimony''; the book, its author warns,
``has been composed not for those who are thorough believers, but for such
as are either wholly unacquainted with the Christian faith, or for those
who, as the apostle terms them, are weak in the faith'' (\work{Against
Celsus}, preface 1--6). What follows is the opposite of weak: point by point, from
Celsus's opening charge that Christians form ``secret associations\ldots\
in violation of the laws'' (\work{Against Celsus} 1.1) to his mockery of a
crucified God, Origen answers with a command of Greek philosophy, history,
and text that no Christian writer had ever displayed and few since have
matched. Eusebius adds that letters survived from Origen ``to the Emperor
Philip, and another to Severa his wife'' (\work{Church History} 6.36.3):
by the 240s the ruined catechist of Alexandria corresponded, at whatever
distance of formality, with the imperial house.

The \work{Against Celsus} is also where Origen's lifelong insistence on
freedom does its most public work: against fate, against astral
determinism, against Celsus's aristocratic contempt for a faith of
``uneducated'' people, he argues that rational creatures are responsible,
that providence teaches rather than compels, and that the gospel's power to
reform the simple is evidence for it, not against it. The book was finished
within about two years of the empire-wide test of exactly those claims.
